% GENERATED FILE. Edit canonical lessons, then run npm run generate:books.
% canonical-source-hash: fnv1a64:6b641cb57edca201
% canonical-lessons: LA-C82-mensa, LA-C82-sella, LA-C82-lectus, LA-C82-fenestra, LA-C82-cubiculum

\chapter{Table, Chair, Bed, Window}
\label{ch:la-table-chair-bed-window}
\hlchaptermodality{82}

\begin{chapteropening}
\textbf{By the end of this chapter:} \emph{I can name a table, a chair, a bed, a window and a bedroom.}

Complete the last lesson of chapter 82: I can name a table, a chair, a bed, a window and a bedroom.
\end{chapteropening}

\section[mēnsa, mēnsae]{mēnsa, mēnsae — a table}

\label{lesson:LA-C82-mensa}



\begin{quote}
Before the new one: say the Latin for a fig, then the Latin for a nut.
\end{quote}

\subsection*{You'll want to know: mēnsa, mēnsae}
\textbf{mēnsa, mēnsae} — \textquotedblleft{}a table\textquotedblright{}.

\begin{cousinweb}
\textbf{mēnsa}, \textquotedblleft{}table\textquotedblright{}, has no sure origin. It is often linked to \emph{mēnsus}, \textquotedblleft{}measured\textquotedblright{}, the participle of \emph{mētior}, \textquotedblleft{}to measure\textquotedblright{}, but the link is not settled.
\end{cousinweb}

\begin{grammarlens}[title={What you've built}]
One more word for this chapter.
\end{grammarlens}

\subsection*{Your turn}
\begin{itemize}
\raggedright
  \item \emph{Say it:} \emph{mēnsa}
  \item \emph{Say it:} \emph{mēnsa}, once more
  \item \emph{Say it:} say \emph{mēnsa}
  \item \emph{Recall:} say the Latin for a fig, then the Latin for a nut, then say \emph{mēnsa} again
\end{itemize}

\begin{tcolorbox}[breakable,colback=teal!4,colframe=teal!35!black,title={Before you move on}]
What does mēnsa, mēnsae mean? (\textquotedblleft{}A table\textquotedblright{}.) Say it once more.
\end{tcolorbox}
\section[sella, sellae]{sella, sellae — a chair, a seat}

\label{lesson:LA-C82-sella}



\begin{quote}
Before the new one: say the Latin for a nut, then the Latin for a table.
\end{quote}

\subsection*{You'll want to know: sella, sellae}
\textbf{sella, sellae} — \textquotedblleft{}a chair, a seat\textquotedblright{}.

\begin{cousinweb}
\textbf{sella} comes from an older \emph{*sedla} (Lewis \& Short), built on Proto-Indo-European \emph{*sed-}, \textquotedblleft{}to sit\textquotedblright{} (EtymDB 2.1), the root of \emph{sedeō}, \textquotedblleft{}to sit\textquotedblright{}. English \emph{saddle} is a cousin from the same root.
\end{cousinweb}

\begin{grammarlens}[title={What you've built}]
One more word for this chapter.
\end{grammarlens}

\subsection*{Your turn}
\begin{itemize}
\raggedright
  \item \emph{Say it:} \emph{sella}
  \item \emph{Say it:} \emph{sella}, once more
  \item \emph{Say it:} say \emph{sella}
  \item \emph{Recall:} say the Latin for a nut, then the Latin for a table, then say \emph{sella} again
\end{itemize}

\begin{tcolorbox}[breakable,colback=teal!4,colframe=teal!35!black,title={Before you move on}]
What does sella, sellae mean? (\textquotedblleft{}A chair, a seat\textquotedblright{}.) Say it once more.
\end{tcolorbox}
\section[lectus, lectī]{lectus, lectī — a bed, a couch}

\label{lesson:LA-C82-lectus}



\begin{quote}
Before the new one: say the Latin for a table, then the Latin for a chair.
\end{quote}

\subsection*{You'll want to know: lectus, lectī}
\textbf{lectus, lectī} — \textquotedblleft{}a bed, a couch\textquotedblright{}.

\begin{cousinweb}
\textbf{lectus} is inherited, usually traced to Proto-Indo-European \emph{*legh-}, \textquotedblleft{}to lie\textquotedblright{} (EtymDB 2.1); Greek \emph{lekhos}, \textquotedblleft{}bed\textquotedblright{}, is a cousin (Lewis \& Short), and so is English \emph{lie}. It is a different word from \emph{lēctus}, \textquotedblleft{}chosen\textquotedblright{}.
\end{cousinweb}

\begin{grammarlens}[title={What you've built}]
One more word for this chapter.
\end{grammarlens}

\subsection*{Your turn}
\begin{itemize}
\raggedright
  \item \emph{Say it:} \emph{lectus}
  \item \emph{Say it:} \emph{lectus}, once more
  \item \emph{Say it:} say \emph{lectus}
  \item \emph{Recall:} say the Latin for a table, then the Latin for a chair, then say \emph{lectus} again
\end{itemize}

\begin{tcolorbox}[breakable,colback=teal!4,colframe=teal!35!black,title={Before you move on}]
What does lectus, lectī mean? (\textquotedblleft{}A bed, a couch\textquotedblright{}.) Say it once more.
\end{tcolorbox}
\section[fenestra, fenestrae]{fenestra, fenestrae — a window}

\label{lesson:LA-C82-fenestra}



\begin{quote}
Before the new one: say the Latin for a chair, then the Latin for a bed.
\end{quote}

\subsection*{You'll want to know: fenestra, fenestrae}
\textbf{fenestra, fenestrae} — \textquotedblleft{}a window\textquotedblright{}.

\begin{cousinweb}
\textbf{fenestra} has no sure origin; it is often thought to be a loan from Etruscan, and the old link with Greek \emph{phainō}, \textquotedblleft{}to show\textquotedblright{}, that Lewis \& Short give is not generally accepted. English \emph{defenestrate} is built on it.
\end{cousinweb}

\begin{grammarlens}[title={What you've built}]
One more word for this chapter.
\end{grammarlens}

\subsection*{Your turn}
\begin{itemize}
\raggedright
  \item \emph{Say it:} \emph{fenestra}
  \item \emph{Say it:} \emph{fenestra}, once more
  \item \emph{Say it:} say \emph{fenestra}
  \item \emph{Recall:} say the Latin for a chair, then the Latin for a bed, then say \emph{fenestra} again
\end{itemize}

\begin{tcolorbox}[breakable,colback=teal!4,colframe=teal!35!black,title={Before you move on}]
What does fenestra, fenestrae mean? (\textquotedblleft{}A window\textquotedblright{}.) Say it once more.
\end{tcolorbox}
\section[cubiculum, cubiculī]{cubiculum, cubiculī — a bedroom}

\label{lesson:LA-C82-cubiculum}



\begin{quote}
Before the new one: say the Latin for a bed, then the Latin for a window.
\end{quote}

\subsection*{You'll want to know: cubiculum, cubiculī}
\textbf{cubiculum, cubiculī} — \textquotedblleft{}a bedroom\textquotedblright{}. English \emph{cubicle}.

\begin{cousinweb}
\textbf{cubiculum} is built on \emph{cubō}, \textquotedblleft{}to lie down\textquotedblright{} (Lewis \& Short), with \emph{-culum}, an ending for a place or tool: a room to lie down in.
\end{cousinweb}

\begin{grammarlens}[title={What you've built}]
One more word for this chapter.
\end{grammarlens}

\subsection*{Your turn}
\begin{itemize}
\raggedright
  \item \emph{Say it:} \emph{cubiculum}
  \item \emph{Say it:} \emph{cubiculum}, once more
  \item \emph{Say it:} say \emph{cubiculum}
  \item \emph{Recall:} say the Latin for a bed, then the Latin for a window, then say \emph{cubiculum} again
\end{itemize}

\begin{tcolorbox}[breakable,colback=teal!4,colframe=teal!35!black,title={Before you move on}]
What does cubiculum, cubiculī mean? (\textquotedblleft{}A bedroom\textquotedblright{}.) Say it once more.
\end{tcolorbox}
